\begin{abstract}
Since the wide adoption of Large Language Models (LLMs), Retrieval Augmented Generation (RAG) has thrived in the specialised literature as a way to ground LLMs' responses on custom knowledge bases. In this context, a central question about the internal mechanics of the model remains unanswered: when new, retrieved evidence becomes sufficient to answer a question, does an LLM undergo an internal state transition that predicts whether it will use that evidence, before it generates a token? Here, we directly investigate this question. To do so, we develop a framework consisting in paired forward passes over semantically defined \emph{evidence trajectories} for multi-hop questions: for each question we present nested contexts that add required paragraphs to answer the question, paired with various forms of distractors. By fitting linear probes on the paired activation update from the internal residual streams of the model, we then show that information about \emph{sufficiency} (the increment that correctly completes the evidence) is recoverable at around $0.95$ AUROC across three model families. We then show a second, orthogonal direction predicting \emph{uptake}, i.e. whether decisive evidence will actually be integrated rather than ignored, at $0.78$ AUROC. Causal tests highlight how the found directions correlate with but do not cause a model behaviour, while a causal direction acting as an abstention gate (whether the model should answer or wait for more evidence) can be found via the mean-differences of sufficient versus insufficient states.
% A text-based baseline, then, shows how \emph{sufficiency} is recoverable from text alone, while \emph{uptake} can be recovered only from internal states. Finally, we perform causal tests highlighting how the discriminative probes find directions that correlate but do not cause the models behaviour, while a causal direction acting as an abstention gate (whether the model should answer or wait for more evidence) can be found via the mean-differences of sufficient versus insufficient states. 
Finally,  we replicate every finding on HotpotQA, on RAGBench's naturally retrieved documents across five domains, and on a synthetic benchmark with paraphrase, partial and contradictory evidence, showing that directions learned on Wikipedia do transfer, but only partially to other domains.
\end{abstract}
